\begin{proof}[Proof of \cref{obj:lem:common-marginal-recipe-transfer}]
Let
\[
\Pi_b:=\text{the branch-\(b\) prior induced by \(\mathcal R\)},\qquad
\theta_b:=\int r_b(w)\,d\Gamma(w),\qquad b\in\{0,1\},
\]
where \(\Gamma\) is the atom-generator measure and \(r_b\) is the raw target of the recipe.  Thus
\[
\Delta=|\theta_1-\theta_0|.
\]
Also put
\[
u:=65536\,n,\qquad v:=65536\,(n+m),
\]
and, for every law \(P\), define the totalized nonnegative recipe scale
\[
s_+(P):=\max\{s_{\mathcal R}(P),0\}.
\]

\begin{enumerate}
\item The witness bounds imply
\[
|\theta_0|\le 1,\qquad |\theta_1|\le 1.
\]
Indeed, each \(\theta_b\) is the \(\Gamma\)-mean of the raw target, the raw target is \(S(w)\tau(P_b(w))\), \(|\tau(P_b(w))|\le1\) on \(\mathcal M_{d,\epsilon}\), and the recipe bounds give \(\int S\,d\Gamma=1\).  Hence
\[
0\le\Delta\le 2.
\]
The intensities satisfy \(u\ge0\) and \(v\ge0\).  The branch priors are probability measures, and the support part of the witness gives the required nonempty model class. % lean: CausalSmith.Stat.SemisupervisedDiscreteAteAnnotationFrontier.commonMarginal_rawPriorCenter_abs_le_one

\item Let \(\mathscr K_{\rm raw}\) be the ordered raw Poisson-count experiment which, at law \(P\), draws independent ordered Poisson samples from the labeled mark law and the auxiliary mark law with intensities
\[
u\,s_+(P)
\qquad\text{and}\qquad
v\,s_+(P).
\]
Let \(\mathsf Q_b\) be the prior-predictive law obtained from \(\Pi_b\) through \(\mathscr K_{\rm raw}\).  We now identify this law as a Markov image of the raw count law in the statement.  Let \(\mathscr K_{\rm cnt}\) denote the recipe's raw Poisson-count kernel before ordering.  Let \(\operatorname{reindex}\) be the deterministic map that rewrites a raw count table as labeled-mark and auxiliary-mark count functions, and let \(\operatorname{order}\) be the paired-histogram kernel that, conditional on those two count functions, orders the corresponding multiset of labeled marks and the corresponding multiset of auxiliary marks uniformly and independently.  Put
\[
\operatorname{Rec}:=\operatorname{order}\circ\operatorname{reindex} .
\]
For every law \(P\), the raw count kernel factors through this reconstruction as
\[
\operatorname{Rec}\circ \mathscr K_{\rm cnt}(P)
=
\mathscr K_{\rm raw}(P).
\]
Indeed, after the deterministic reindexing, the labeled and auxiliary count functions have independent Poisson coordinates with means \(u\,s_+(P)\) times the labeled mark law and \(v\,s_+(P)\) times the auxiliary mark law; ordering each histogram gives exactly the ordered independent Poisson samples with those two intensities.  Therefore, for each branch,
\[
\mathcal R_{\rm raw,b}
:=\int \mathscr K_{\rm cnt}(P)\,d\Pi_b(P)
\qquad\text{satisfies}\qquad
\mathsf Q_b
=\int \mathscr K_{\rm raw}(P)\,d\Pi_b(P)
=(\operatorname{Rec})_{\#}\mathcal R_{\rm raw,b}.
\]
Total variation contracts under a Markov kernel, and the assumed raw mixture affinity gives
\[
\operatorname{TV}(\mathsf Q_0,\mathsf Q_1)
\le
\operatorname{TV}(\mathcal R_{\rm raw,0},\mathcal R_{\rm raw,1})
\le \frac18 .
\] % lean: CausalSmith.Stat.SemisupervisedDiscreteAteAnnotationFrontier.commonMarginal_predictive_reconstruction_contraction

\item Set
\[
\alpha_{\rm floor}:=\frac1{100}.
\]
The fixed experiment has the parametric two-point floor
\[
\frac{\alpha_{\rm floor}}{n}\le R_\epsilon(n,m,d).
\]
To see this, take
\[
\delta_n:=\frac{2}{5\sqrt n}.
\]
Choose a cell \(x_\star\in[d]\).  Under both laws below, let \(X=x_\star\) almost surely and let the treatment propensity be \(1/2\).  Let the control outcome mean be \(1/2\) under both laws; let the treated outcome mean be \(1/2\) for \(P_{\rm par,0}\) and \(1/2+\delta_n\) for \(P_{\rm par,1}\).  Since \(n\ge1\),
\[
0\le\delta_n\le\frac25,
\]
so these are valid binary outcome laws, and since \(0<\epsilon<1/2\), both belong to \(\mathcal M_{d,\epsilon}\).  Their treatment-covariate marginals agree, their average treatment effects are
\[
\tau(P_{\rm par,0})=0,
\qquad
\tau(P_{\rm par,1})=\delta_n,
\]
and hence the auxiliary samples have identical laws.  For the labeled samples, the standard Bernoulli product testing bound applies because
\[
n\cdot
\frac{(1/2)\delta_n^2}{(1/2)(1-1/2)}
=
\frac{8}{25}
\le \log 2,
\]
and therefore
\[
\operatorname{TV}\!\bigl(P_{\rm par,0}^{\otimes n},P_{\rm par,1}^{\otimes n}\bigr)
\le\frac12.
\]
Le Cam's two-point bound with separation \(\delta_n\) yields
\[
R_\epsilon(n,m,d)
\ge
\frac{(\delta_n/2)^2}{4}
=
\frac1{100\,n}.
\]
For the fixed-sample decision problem with the class-gated clipped squared loss, let
\[
\mathfrak R_{\rm dec}
:=
\inf_{\psi}
\sup_P
\int
\ell_\epsilon(P,\psi)\,d\mathscr K_{\rm fix}(P),
\]
where the infimum is over measurable fixed-sample rules, \(\mathscr K_{\rm fix}\) is the fixed labeled-and-auxiliary experiment, and
\[
\ell_\epsilon(P,z)
:=
\begin{cases}
\bigl(\operatorname{clip}_{[-1,1]}(z)-\tau(P)\bigr)^2,& P\in\mathcal M_{d,\epsilon},\\
0,& P\notin\mathcal M_{d,\epsilon}.
\end{cases}
\]
We use the following elementary comparison between \(\cref{obj:def:minimax-risk}\) and the class-gated decision problem:
\[
R_\epsilon(n,m,d)
\le
\mathfrak R_{\rm dec}.
\]
Indeed, every measurable fixed-sample rule \(\psi\) induces the measurable estimator
\[
\widetilde T_\psi:=\operatorname{clip}_{[-1,1]}\circ\psi .
\]
For every \(P\in\mathcal M_{d,\epsilon}\), the class-gated loss evaluates the same clipped rule, so
\[
\int \ell_\epsilon(P,\psi(y))\,d\mathscr K_{\rm fix}(P)(y)
=
E_{P^{\otimes n}\otimes P_{XA}^{\otimes m}}
\!\left[\bigl(\widetilde T_\psi-\tau(P)\bigr)^2\right].
\]
Thus the worst-case on-class decision risk of \(\psi\) is exactly the worst-case mean-squared risk of \(\widetilde T_\psi\).  Since \(\widetilde T_\psi\) is one admissible estimator in \(\mathcal T_{n,m,d}\), \(\cref{obj:def:minimax-risk}\) is no larger than that worst-case risk; taking the infimum over \(\psi\) gives the displayed comparison.  Combining it with the parametric floor gives
\[
\frac{\alpha_{\rm floor}}{n}
\le
\mathfrak R_{\rm dec}.
\] % lean: CausalSmith.Stat.SemisupervisedDiscreteAteAnnotationFrontier.commonMarginal_explicit_parametric_floor

\item If \(\Delta=0\), then the desired lower bound is \(0\le R_\epsilon(n,m,d)\), which follows from the preceding floor since \(n\ge1\).  Hence assume for the rest of the proof that
\[
0<\Delta.
\]
The variance-scale hypothesis and \(C_\epsilon>0\) give
\[
C_\epsilon kBa\le \frac{\Delta^2}{65536}.
\]
By \cref{obj:lem:common-marginal-recipe-target-concentration}, for each \(b\in\{0,1\}\),
\[
\Pi_b\!\left\{P:\frac{\Delta}{4}<|\tau(P)-\theta_b|\right\}
\le
\frac{512\,C_\epsilon kBa}{\Delta^2}
\le
\frac1{128}
\le
\frac1{16}.
\]
The assumed common-marginal Poisson-prior bounds also include the raw-center separation
\[
\delta_\epsilon k a\le |\theta_1-\theta_0|=\Delta,
\]
which is the center-separation hypothesis required below.
By \cref{obj:lem:common-marginal-recipe-scale-tail}, again for each \(b\in\{0,1\}\),
\[
\Pi_b\{P:s_+(P)<1/2\}
\le
4C_\epsilon kBa
\le
\frac{\Delta^2}{16384}.
\] % lean: CausalSmith.Stat.SemisupervisedDiscreteAteAnnotationFrontier.commonMarginal_prior_target_concentration

\item The two fuzzy hypotheses with centers \(\theta_0,\theta_1\), radius \(\Delta/4\), prior tail bounds \(1/16,1/16\), and predictive total variation bound \(1/8\) give the raw lower bound as follows.  For a measurable raw rule \(\varphi\), define
\[
\mathcal B_b(\varphi)
:=
\int^{+}
\int^{+}
\ell_\epsilon(P,\varphi(y))\,d\mathscr K_{\rm raw}(P)(y)\,d\Pi_b(P)
\in\overline{\mathbb R}_+,
\qquad b\in\{0,1\}.
\]
These are the branch Bayes decision risks for the class-gated clipped loss.  Since the two branch priors are supported on \(\mathcal M_{d,\epsilon}\), the integrand on the support is
\[
\ell_\epsilon(P,\varphi(y))
=
\bigl(\operatorname{clip}_{[-1,1]}(\varphi(y))-\tau(P)\bigr)^2.
\]
Let \(\mathsf Q_b\) be the branch predictive laws from step 2, and first suppose \(\theta_0\le\theta_1\).  Put
\[
\theta_{\rm mid}:=\frac{\theta_0+\theta_1}{2},
\qquad
\mathsf A_\varphi
:=
\bigl\{y:\operatorname{clip}_{[-1,1]}(\varphi(y))\ge\theta_{\rm mid}\bigr\},
\]
and define the good prior events
\[
\mathsf G_0:=\bigl\{P:|\tau(P)-\theta_0|\le\Delta/4\bigr\},
\qquad
\mathsf G_1:=\bigl\{P:|\tau(P)-\theta_1|\le\Delta/4\bigr\}.
\]
Because \(\theta_1-\theta_0=|\theta_1-\theta_0|\ge\Delta\), on \(\mathsf G_0\cap\mathsf A_\varphi\) the clipped decision is at least \(\Delta/4\) above \(\tau(P)\), while on \(\mathsf G_1\cap\mathsf A_\varphi^c\) it is at least \(\Delta/4\) below \(\tau(P)\).  Hence
\[
\mathcal B_0(\varphi)
\ge
\frac{\Delta^2}{16}
\Pr_{\Pi_0\otimes\mathscr K_{\rm raw}}(\mathsf G_0\cap\mathsf A_\varphi),
\qquad
\mathcal B_1(\varphi)
\ge
\frac{\Delta^2}{16}
\Pr_{\Pi_1\otimes\mathscr K_{\rm raw}}(\mathsf G_1\cap\mathsf A_\varphi^c).
\]
The concentration bounds from step 4 give
\[
\Pr_{\Pi_0\otimes\mathscr K_{\rm raw}}(\mathsf G_0^c)\le\frac1{16},
\qquad
\Pr_{\Pi_1\otimes\mathscr K_{\rm raw}}(\mathsf G_1^c)\le\frac1{16},
\]
and therefore
\[
\mathcal B_0(\varphi)
\ge
\frac{\Delta^2}{16}\left(\mathsf Q_0(\mathsf A_\varphi)-\frac1{16}\right),
\qquad
\mathcal B_1(\varphi)
\ge
\frac{\Delta^2}{16}\left(1-\mathsf Q_1(\mathsf A_\varphi)-\frac1{16}\right).
\]
Using \(\max\{x_0,x_1\}\ge(x_0+x_1)/2\) and the total-variation bound,
\[
\begin{aligned}
\max\{\mathcal B_0(\varphi),\mathcal B_1(\varphi)\}
&\ge
\frac{\Delta^2}{32}
\left(1-\bigl(\mathsf Q_1(\mathsf A_\varphi)-\mathsf Q_0(\mathsf A_\varphi)\bigr)-\frac1{16}-\frac1{16}\right)\\
&\ge
\frac{\Delta^2}{32}
\left(1-\frac18-\frac1{16}-\frac1{16}\right)
=
\frac{3\Delta^2}{128}.
\end{aligned}
\]
If \(\theta_1\le\theta_0\), the same calculation is applied with the two branches interchanged.  Thus, for every measurable raw rule \(\varphi\),
\[
\frac{3\Delta^2}{128}
\le
\max\Bigl\{
\mathcal B_0(\varphi),\mathcal B_1(\varphi)
\Bigr\}
\]
in \(\overline{\mathbb R}_+\). % lean: CausalSmith.Stat.SemisupervisedDiscreteAteAnnotationFrontier.commonMarginal_raw_fuzzy_lower

\item The loss is bounded by
\[
0\le \ell_\epsilon(P,z)\le 4
\]
because both \(\operatorname{clip}_{[-1,1]}(z)\) and \(\tau(P)\) lie in \([-1,1]\) on the model class, and the loss is \(0\) off the model class.  The Poisson retention tails used to transfer raw rules to fixed-sample rules are as follows.  For any integer \(\mathfrak n\ge1\) and any Poisson mean \(\lambda\ge32768\,\mathfrak n\),
\[
\Pr\{\operatorname{Poi}(\lambda)<\mathfrak n\}
\le
\exp(-\lambda/8)
\le
\frac1{4096\,\mathfrak n}.
\]
Indeed, with \(z_\lambda=\lambda/8\),
\[
\sqrt{2\lambda z_\lambda}=\lambda/2,
\]
and \(\{\operatorname{Poi}(\lambda)<\mathfrak n\}\subseteq
\{\lambda-\operatorname{Poi}(\lambda)>\lambda/2\}\); the Poisson Bernstein lower-tail bound gives the exponential term, and \(\exp(-z_\lambda)\le 1/z_\lambda\le1/(4096\mathfrak n)\).

On the event \(s_+(P)\ge1/2\), the labeled Poisson mean satisfies
\[
u\,s_+(P)\ge 32768\,n,
\]
so the labeled retention failure probability is at most \(1/(4096n)\).  Similarly,
\[
v\,s_+(P)\ge 32768\,(n+m),
\]
and the event of retaining fewer than \(m\) auxiliary observations is contained in the event of retaining fewer than \(n+m\) observations; hence the auxiliary retention failure probability is at most
\[
\frac1{4096(n+m)}\le \frac1{4096n}.
\] % lean: CausalSmith.Stat.SemisupervisedDiscreteAteAnnotationFrontier.commonMarginal_poisson_lower_tail

\item We next transfer the raw fuzzy lower bound to the fixed experiment by the prefix coupling.  Fix any measurable fixed-sample decision rule \(\psi\).  Define its transferred raw rule \(\varphi_\psi\) as follows.  Given an ordered raw output \(y_{\rm raw}\), set
\[
\varphi_\psi(y_{\rm raw})
=
\begin{cases}
\psi\bigl(\text{the first \(n\) labeled marks and the first \(m\) auxiliary marks in \(y_{\rm raw}\)}\bigr),
& \text{if those marks are present},\\
\zeta_\star,& \text{otherwise}.
\end{cases}
\]
The rule uses only the first \(m\) auxiliary marks, while the raw auxiliary Poisson process has intensity \(v\,s_+(P)=65536(n+m)s_+(P)\); the auxiliary tail estimate below therefore bounds the event of having fewer than \(m\) auxiliary marks by the larger-threshold event of having fewer than \(n+m\).  Here \(\zeta_\star\in\mathbb R\) is any fixed real-valued action.  For a law \(P\), let
\[
\mathsf F_P
:=
\bigl\{\text{the ordered raw draw contains fewer than \(n\) labeled marks or fewer than \(m\) auxiliary marks}\bigr\}.
\]
On \(\mathsf F_P^c\), the first \(n\) labeled marks and first \(m\) auxiliary marks have exactly the fixed-sample law \(\mathscr K_{\rm fix}(P)\), because the ordered raw experiment is generated by independent Poisson samples from the same labeled and auxiliary mark kernels and the fixed experiment keeps the first prescribed counts.

For branch \(b\), the exceptional probability under the branch prior predictive is bounded by integrating the conditional failure probability.  Splitting according to the scale event gives
\[
\begin{aligned}
\Pr_{\Pi_b\otimes\mathscr K_{\rm raw}}(\mathsf F_P)
&\le
\Pi_b\{P:s_+(P)<1/2\} \\
&\quad+
\int_{\{P:s_+(P)\ge1/2\}}
\left[
\Pr\{\operatorname{Poi}(u s_+(P))<n\}
+
\Pr\{\operatorname{Poi}(v s_+(P))<m\}
\right]d\Pi_b(P).
\end{aligned}
\]
The scale-tail bound from step 4 and the Poisson retention bounds from step 6 therefore yield, for each \(b\in\{0,1\}\),
\[
\Pr_{\Pi_b\otimes\mathscr K_{\rm raw}}(\mathsf F_P)
\le
\frac{\Delta^2}{16384}
+\frac1{4096n}
+\frac1{4096n}.
\]
Since \(0\le\ell_\epsilon\le4\), the Bayes risk of the transferred raw rule satisfies
\[
\begin{aligned}
\mathcal B_b(\varphi_\psi)
&\le
\int
\int \ell_\epsilon(P,\psi(y))\,d\mathscr K_{\rm fix}(P)(y)\,d\Pi_b(P)\\
&\quad+
4\left(
\frac{\Delta^2}{16384}
+\frac1{4096n}
+\frac1{4096n}
\right)\\
&\le
\sup_{P}\int \ell_\epsilon(P,\psi(y))\,d\mathscr K_{\rm fix}(P)(y)
+
4\left(
\frac{\Delta^2}{16384}
+\frac1{4096n}
+\frac1{4096n}
\right).
\end{aligned}
\]
Combining this with the raw fuzzy bound above gives, for every measurable \(\psi\),
\[
\frac{3\Delta^2}{128}
\le
\sup_{P}\int \ell_\epsilon(P,\psi(y))\,d\mathscr K_{\rm fix}(P)(y)
+
4\left(
\frac{\Delta^2}{16384}
+\frac1{4096n}
+\frac1{4096n}
\right).
\]
Equivalently,
\[
\frac{3\Delta^2}{128}
-
4\left(
\frac{\Delta^2}{16384}
+\frac1{4096n}
+\frac1{4096n}
\right)
\le
\sup_{P}\int \ell_\epsilon(P,\psi(y))\,d\mathscr K_{\rm fix}(P)(y).
\]
Taking the infimum over \(\psi\) gives
\[
\frac{3\Delta^2}{128}
-
4\left(
\frac{\Delta^2}{16384}
+\frac1{4096n}
+\frac1{4096n}
\right)
\le
\mathfrak R_{\rm dec}.
\]
The floor inequality from step 3 also gives \(\alpha_{\rm floor}/n\le\mathfrak R_{\rm dec}\), so
\[
\max\left\{
\frac{3\Delta^2}{128}
-
4\left(
\frac{\Delta^2}{16384}
+\frac1{4096n}
+\frac1{4096n}
\right),
\,
\frac{\alpha_{\rm floor}}{n}
\right\}
\le
\mathfrak R_{\rm dec}.
\]
For the reverse comparison to step 3, start with any measurable estimator \(T\) in
\cref{obj:def:minimax-risk} and use the same function as the fixed-sample decision rule
\(\psi_T\).  On \(\mathcal M_{d,\epsilon}\), clipping cannot increase squared error because
\(\tau(P)\in[-1,1]\), while off the model class the class-gated loss is zero.  Hence the
worst-case decision risk of \(\psi_T\) is no larger than the worst-case mean-squared risk of
\(T\).  Taking the infimum over \(T\) gives
\[
\mathfrak R_{\rm dec}
\le
R_\epsilon(n,m,d),
\]
and consequently
\[
\max\left\{
\frac{3\Delta^2}{128}
-
4\left(
\frac{\Delta^2}{16384}
+\frac1{4096n}
+\frac1{4096n}
\right),
\,
\frac{\alpha_{\rm floor}}{n}
\right\}
\le
R_\epsilon(n,m,d).
\] % lean: Causalean.Stat.FiniteRaoBlackwell.IndependentPoissonPrefix.RandomScaleMinimaxTransfer.randomScale_twoFuzzy_minimax_lower_transfer

\item It remains only to read off the deterministic arithmetic.  From the floor bound,
\[
R_\epsilon(n,m,d)\ge \frac{\alpha_{\rm floor}}{n}>0.
\]

First suppose
\[
\frac1n\le \Delta^2.
\]
Then
\[
4\left(\frac{\Delta^2}{16384}
+\frac1{4096n}
+\frac1{4096n}\right)
\le
4\left(\frac{\Delta^2}{16384}
+\frac{2\Delta^2}{4096}\right)
=
\frac{9\Delta^2}{4096}.
\]
Therefore
\[
\frac{3\Delta^2}{128}
-
4\left(\frac{\Delta^2}{16384}
+\frac1{4096n}
+\frac1{4096n}\right)
\ge
\frac{87\Delta^2}{4096}
\ge
\frac{\Delta^2}{65536}.
\]
The first branch of the maximum in the preceding display gives
\[
R_\epsilon(n,m,d)\ge \frac{\Delta^2}{65536}.
\]

Now suppose
\[
\Delta^2<\frac1n.
\]
Since \(\alpha_{\rm floor}=1/100\),
\[
\frac{\Delta^2}{65536}
<
\frac1{65536\,n}
\le
\frac1{100\,n}
=
\frac{\alpha_{\rm floor}}{n}.
\]
The second branch of the maximum gives the same conclusion.  Combining the two cases,
\[
R_\epsilon(n,m,d)\ge \frac{\Delta^2}{65536}.
\] % lean: CausalSmith.Stat.SemisupervisedDiscreteAteAnnotationFrontier.fixed_sample_common_marginal_recipe_transfer_rule
\end{enumerate}
\end{proof}
